\section{Discussion}\label{sec:discussion}
\paragraph{What the simulations say.}
\begin{itemize}
\item \textbf{The two best published electrochemical specifications can train.} Both the Li-ECRAM fit and the protonic device
      reach the exact-update baseline within our gap lines, or within the statistical resolution of five seeds, on all three
      tasks, provided a Tiki-Taka-family algorithm is used.
\item \textbf{Plain SGD never does,} on any electrochemical device we tried.
\item \textbf{Width is not the obstacle} we feared. Once each arm has its own learning rate, the gap does not grow with width
      over the factor of four we could afford.
\item \textbf{Two device properties decide.} The decisive specifications are not the number of levels, which TTv2 needs only in
      the hundreds, but:
      \begin{itemize}
      \item an accurately referenced symmetry point, or an algorithm that does not need one;
      \item a saturating (asymmetric) response that gives the gradient accumulator its decay.
      \end{itemize}
      Retention of several training lengths and write noise of up to one or two steps are tolerated.
\item \textbf{The energy case rests on the periphery,} not on the device. Femtojoule pulses make the update almost free. The
      forward and backward reads, priced at a measured analog-AI chip's efficiency, set the system at about
      \REnRatioMin$\times$ an H100's energy per weight-row: useful, but not the thousandfold figures that device-only accounting
      suggests.
\end{itemize}

\paragraph{Widrow's memistor.} Taken at its inventor's description (thousands of near-linear levels), the memistor would train
as well as modern ECRAM on our tasks. The obstacles of 1960 were the missing multilayer rule, wet macroscopic cells and an
energy per level of tens of microjoules. Accuracy was not the obstacle. The idea survived; the implementation had to shrink by
about \ROrdersWidrowVsProtonic{} orders of magnitude in energy per update.

\subsection{Limitations}\label{sec:limits}
\begin{itemize}
\item \textbf{Simulation only.} Every device is a parametric model.
      \begin{itemize}
      \item \RNCatalogEstimated{} of the \RNCatalog{} catalog entries rest partly or wholly on estimates.
      \item The protonic device's write noise and spreads are estimates. A 2026 measurement on BEOL-compatible protonic
            ECRAM~\cite{xu2026} reports 3\% cycle-to-cycle and 21\% device-to-device variability. These lie inside our tolerated
            region, but we did not refit the catalog to them.
      \item The model has no state-dependent noise beyond the soft bounds, no history effects or relaxation transients
            beyond a single drift time constant, no temperature, and no endurance limit.
      \end{itemize}
\item \textbf{Ideal reads.} We isolate the update. Read noise, IR drop, ADC/DAC quantisation and output bounds would add to
      every gap; our validation omits them and still matches the published numbers. Real arrays also have line resistance and
      sneak paths. Three-terminal devices avoid the latter but need access transistors.
\item \textbf{Small models and short budgets.}
      \begin{itemize}
      \item The networks have $10^5$ to $10^6$ weights.
      \item CIFAR uses a 10{,}000-image subset whose seed noise (\REThreeCifarFpSd\,pp) limits the resolution at the 2\,pp
            line.
      \item The LM is a two-layer byte model.
      \item The width study spans only a factor of four.
      \end{itemize}
      Larger models are where the remaining questions are. Simulating every pulse becomes expensive; this is also the
      experience reported by~\cite{rasch2024}.
\item \textbf{The baseline is SGD.} The FP baseline and every analog arm are plain-SGD-family learners. Against Adam the best
      analog LM is still about as far behind as FP-SGD is.
\item \textbf{Mini-batches.} Our update applies a mini-batch's summed pulses as interleaved, rather than sample by sample. The
      validation shows that batch 16 reproduces published batch-1 results, but strongly asymmetric devices degrade at larger
      batches.
\item \textbf{Algorithm details.} We implemented TTv1, TTv2, c-TTv2 and AGAD from the papers' text, not from the authors'
      code. Every TT-family arm uses a multi-pulse transfer (up to 31 pulses), and a single settings selection on MNIST is
      shared by all tasks. Results for c-TTv2 and AGAD could improve with per-task tuning.
\item \textbf{Energy.}
      \begin{itemize}
      \item The MVM efficiency comes from phase-change inference chips.
      \item ECRAM gate drivers and access devices have no published system-level energy.
      \item The update periphery is our estimate.
      \item The GPU figure uses board power and a published utilisation; it ignores host, network and cooling, which would
            raise the GPU's energy.
      \end{itemize}
\item \textbf{Statistics.} Five seeds per cell, no multiplicity correction, and single seeds for most points of the
      sensitivity grids. Several cells are UNRESOLVED; more seeds, not more conditions, would settle them.
\end{itemize}
